\section{Short square relations and the rank of a pair}
\label{p13:rig:section}

A large intersection of the square-class spaces of two windows
contains many relations, but those relations need not initially have
short support.  We first extract a large subspace with a short basis.
Two independent basis vectors then produce an algebraic curve through
the two window centres.  A resultant argument puts the whole subspace
on one curve, while the relations force many integer points in the
small grid of offsets.  A point count uniform in the degree rules out
this configuration unless the centres lie on a rational alignment.
The distinction between a fixed degree and a growing degree is
essential: the basis supports grow with $L$ throughout the argument.

The combination of resultant elimination and degree-uniform curve
bounds also appears in Shparlinski and Sleiman's work on consecutive
multiplicatively dependent triples
\cite[Theorem~1.1 and Section~3]{shparlinskisleiman2026}.
Here the common curve must accommodate an entire space of square-class
relations and the integer fibres forced by that space in the offset grid.

All estimates in this section are uniform in
$x,y\in I_N$, $|L-\log_2N|\le C$, for fixed $C$, and we put
$B=L+1$.  A logarithm without a subscript is natural.
The conclusion used in the probability argument is
Corollary~\ref{p13:rig:triple}.

\subsection{Square classes on all windows}
\label{p13:rig:notation}

We use the true window
$V_x=\{x-1,x,\ldots,x+L-1\}$, including its left endpoint.
Let $E_x$ be the space of even selections from $V_x$, and let
\[
 \phi_x:E_x\longrightarrow
 \mathbb Q_{>0}^{\times}/\mathbb Q_{>0}^{\times2},
 \qquad I\longmapsto\left[\prod_{n\in I}n\right]
\]
be the square-class map.  Its image is the space $S_x$ already used
above.  Put $d_{xy}=\dim(S_x\cap S_y)$ and, for any finite set of
centres $T$, define the \emph{relative nullity}
\begin{equation}
 \rho_{\rm rel}(T)=\sum_{x\in T}\dim S_x-
                         \dim\left(\sum_{x\in T}S_x\right).
 \label{p13:rig:relative-nullity}
\end{equation}
On the regular set, where every $S_x$ has dimension $L$, this is
the earlier nullity $\rho(T)$.  Formula
\eqref{p13:rig:relative-nullity} also applies to the remaining
windows and separates their local rank defects from dependence
between distinct windows.

A \emph{complete relation} of a pair is a pair of even selections
with the same square class.  We represent its selections by offsets
$I,J\subset\{-1,0,\ldots,L-1\}$ and attach the polynomial
\begin{equation}
 F_w(X,Y)=\prod_{u\in I}(X+u)\prod_{v\in J}(Y+v),
 \qquad w=(I,J).
 \label{p13:rig:word-polynomial}
\end{equation}
Thus $F_w(x,y)$ is an integer square.  The word $w$ is a vector
over $\F_2$; addition means symmetric difference in each block.
The two blocks remain disjoint formal copies even when the actual
windows overlap.  Consequently, independent words give independent
classes of their polynomials in
$\mathbb Q(X,Y)^\times/\mathbb Q(X,Y)^{\times2}$: the factors
$X+u$ and $Y+v$ are distinct irreducible polynomials.

We call $x,y$ \emph{aligned} if there are coprime positive integers
$a,b\le B^3$ such that
\begin{equation}
                   |bx-ay|\le(a+b)B.
 \label{p13:rig:alignment}
\end{equation}
This definition includes nearby centres.  Otherwise they are
nonaligned.  Alignment is a property of the centres; compatibility
of their run events is a separate condition.


\subsection{Extracting short relations}

Primes larger than $10B$ meet each physical block at most once.  Their
parity equations therefore join positions into components; a component
meeting a prime only once is forced to zero.  Local square-product
separation shows that there are only $O_C(B/\log B)$ defective
positions.  Thus the number of free components meeting both blocks is
$d_{xy}+O_C(B/\log B)$; see
Lemmas~\ref{p13:rig:defective} and~\ref{p13:rig:components}.

For medium primes, columns of weight two can be kept as edges of a
graph.  Its short cycles cancel their equations.  Discarding only
columns of weight at least three costs one third of their total
weight, explaining the factor $2/3$ in the following extraction.

\begin{theorem}[A short basis with a controlled loss of dimension]
\label{p13:rig:short-basis}
For every fixed $a\in(0,1)$ there is a subspace $\mathcal K_a$
of complete relations of the pair, with injective projections
onto each physical block and onto the common square class,
such that
\begin{equation}
 \dim\mathcal K_a\ge d_{xy}-\frac23B\log(1/a)
                    -O_{a,C}\!\left(\frac B{\sqrt{\log B}}\right).
 \label{p13:rig:short-dimension}
\end{equation}
It has a basis whose vectors have total physical support at most
\begin{equation}
                    O_{a,C}\!\left(
                         \frac{B^a\log B}{\log\log B}\right).
 \label{p13:rig:short-support}
\end{equation}
Every nonzero vector selects a nonempty even set in each block
and has nonzero common square class.  Neither compatibility nor
nonalignment is assumed.
\end{theorem}

The full extraction proof, including the final small-prime equations
and projection injectivity, is in
Section~\ref{p14:comp:rigidity:short-basis}.

\subsection{One curve for a whole relation space}

Two independent short words first yield a nonzero polynomial through
the window centres.  Proposition~\ref{p13:rig:two-witnesses} constructs
it as the norm of a biquadratic auxiliary function, with degree and
logarithmic height linear in the support size.  Polynomial gcd removal
is performed before specialization, so vanishing at the integer point
does not destroy the norm argument.

The simultaneous step below is stronger than constructing a separate
curve for each word.  Pairwise resultants have quadratic height cost
in the support, independent of the number of basis vectors.

\begin{theorem}[A common curve]
\label{p13:rig:common-curve}
There is an absolute constant $c>0$ with the following property.
Let $V$ be a space of even formal words of dimension $t\ge2$,
with a basis $w_1,\ldots,w_t$ satisfying
$|I_{w_i}|,|J_{w_i}|\le2q_0$, where $1\le q_0\le B$.
Suppose all $F_{w_i}(x,y)$ are integer squares and
$x,y\in[N,2N)$.  If $N>(B+1)^{cq_0^2}$, there is a
primitive irreducible polynomial $C_0\in\Z[X,Y]$, depending
on both variables, such that
\begin{equation}
 C_0(x,y)=0,\qquad \deg C_0\le48q_0,
 \qquad H(C_0)\le(B+1)^{cq_0}.
 \label{p13:rig:common-degree-height}
\end{equation}
In the function field $K_0=\mathbb Q(C_0=0)$, the image of
\[
               V\longrightarrow K_0^\times/K_0^{\times2},
                  \qquad w\longmapsto[F_w]
\]
has dimension at most one.  The constant is independent of $t$.
\end{theorem}

The proof is given in Section~\ref{p14:comp:rigidity:common-curve}.
Its key point is that every norm polynomial must share the same
irreducible factor through $(x,y)$: otherwise a resultant would give
a nonzero polynomial of height $(B+1)^{O(q_0^2)}$ with integer root
$x\ge N$.  Primitivity of the auxiliary coefficients then forces the
functional square-class image to have dimension at most one.

\subsection{Functional squares force integer fibres}

The centres have size $N$, but the next step counts points
in a box of side $O(B)$.  The bridge is that a square
function has even valuation at every point of the normalized
curve.  A simple zero belonging to a selected first-block
offset must therefore be matched by a selected second-block
offset.

\begin{proposition}[Integer fibres in the offset grid]
\label{p13:rig:grid-fibres}
Under the hypotheses of Theorem~\ref{p13:rig:common-curve},
assume both physical projections of $V$ are injective.
Put $D=\deg C_0$, $d_X=\deg_XC_0$, $d_Y=\deg_YC_0$.
For an absolute constant $c$,
\begin{equation}
 \#\{(u,v)\in[-B,B]^2\cap\Z^2:C_0(-u,-v)=0\}
       \ge\max(d_X,d_Y)(t-1-cD^2).
 \label{p13:rig:grid-count}
\end{equation}
A negative right side may be replaced by zero.
\end{proposition}

\begin{proof}
The kernel $V_0$ of the map to functional square classes
has dimension at least $t-1$.  The leading coefficients
and discriminants of $C_0$ in each variable are nonzero
polynomials, of total degree $O(D^2)$ in the remaining
variable.  Exclude the offsets at which one of these
polynomials vanishes.  Requiring the corresponding
coordinates of a word to be zero leaves a subspace $V_1$
of dimension at least $t-1-cD^2$.

Let $\mathcal U_X$ be the union of first-block offsets used
by $V_1$.  Projection injectivity gives
$|\mathcal U_X|\ge\dim V_1$.
For $u\in\mathcal U_X$, select a word containing $X+u$.
The noncritical fibre $X=-u$ has exactly $d_Y$ distinct
affine roots on the geometric normalization.  At each,
$X+u$ has a simple zero and no other first-block factor
vanishes.  The product for the selected word is a square
function, so a selected factor $Y+v$ must vanish there.
Every root is thus an integer $-v$ in the offset interval.
Distinct $u$ give disjoint fibres, proving at least
$d_Y\dim V_1$ grid points.  Interchanging $X,Y$ proves
the bound with $d_X$ as well.
The argument may be performed after extension to
$\mathbb C$ and normalization of the components; geometric
integrality is not required.
\end{proof}

\subsection{The point count at growing degree}

We use the integer-point part of Binyamini--Cluckers--Novikov
\cite[Theorem~2, equation~(27), in the specified author
version]{binyaminicluckersnovikov2024}.

\begin{theorem}[Binyamini--Cluckers--Novikov]
\label{p13:rig:bcn}
For an integral algebraic curve $\mathcal C\subset\mathbb A^n_K$
of degree $D$, where $n>1$ and $K\subset\mathbb C$, there are
a constant $c(n)$ and an absolute constant $\kappa$ such that
\[
 \#\{z\in\mathcal C\cap\Z^n:\max_i|z_i|\le H\}
       \le c(n)D^2H^{1/D}(\log H)^\kappa,
       \qquad H>2.
\]
The constants are independent of the curve's coefficients
and degree.  Integrality over $K$ is sufficient.
\end{theorem}

For curves over $\mathbb Q$, a corresponding affine estimate is also
given by Binyamini, Cluckers, and Kato
\cite[Corollary~8, equation~(23)]{binyaminicluckerskato2025}.
The external theorem above asserts the existence of an absolute
logarithmic exponent. The admissible value $\kappa=17$ is supplied by
the pinned Lean declaration
\href{https://github.com/Vulkin-prog/paper-p1-lean/blob/a1c8092f7e25acf7f3b78283039cfbd69491f88e/PaperP1/Arithmetic/IntegralAffineCurvePointCount.lean}
{\texttt{article\_13\_4}}
in the namespace \texttt{PaperP1.IntegralAffineCurvePointCount}.
Only the existence of a fixed exponent is used below.

\begin{corollary}[Sublinear counts in the offset box]
\label{p13:rig:sublinear-grid}
Fix $0\le a<1/2$ and $c_0>0$. Let $D_{\max}=D_{\max}(B)$
be any fixed degree envelope such that, for every $\delta>0$,
$D_{\max}(B)\le B^{a+\delta}$ for all sufficiently large $B$.
Uniformly for integral plane curves of degree $2\le D\le D_{\max}(B)$,
\[
 \#(\mathcal C\cap[-c_0B,c_0B]^2\cap\Z^2)
 \ll_{c_0}(B^{1/2}+D_{\max}^2)(\log B)^\kappa=o(B).
\]
\end{corollary}

\begin{proof}
Set $H=\max(3,c_0B)$.  The logarithmic derivative of
$u^2H^{1/u}$ is $(2u-\log H)/u^2$, so its maximum on
$[2,M]$, $M\ge2$, is attained at an endpoint.
Put $M'=\max(M,\log H)$. Applying this observation on $[2,M']$,
and using $H^{1/M'}\le e$, gives, for every $2\le D\le M$,
\[
 D^2H^{1/D}
 \le 4H^{1/2}+e(M')^2
 \le4H^{1/2}+e\{M^2+(\log H)^2\}.
\]
The last logarithmic square is absorbed into $B^{1/2}$.
Apply Theorem~\ref{p13:rig:bcn} with $n=2$ and $M=D_{\max}(B)$.
Choose a fixed $\delta>0$ with $2(a+\delta)<1$. Both
$B^{1/2}(\log B)^\kappa$ and
$B^{2(a+\delta)}(\log B)^\kappa$ are $o(B)$.
The threshold is chosen after $c_0,a$, the envelope and the requested
error, and before the curve and its coefficients. No lower growth of
$D_{\max}$ is needed; bounded and polylogarithmic envelopes are included.
\end{proof}

This application uses the theorem's degree-uniform constant.
A bound with an unspecified constant depending on $D$ would
not give the conclusion of Corollary~\ref{p13:rig:sublinear-grid}.

\subsection{Nonaligned pairs have rank less than half a window}

\begin{theorem}[A uniform nonaligned rank bound]
\label{p13:rig:nonaligned}
For every fixed critical band,
\begin{equation}
 \limsup_{N\to\infty}\ 
 \sup_{\substack{x,y\in I_N\\x,y\ {\rm nonaligned}}}
             \frac{d_{xy}}B
       \le\frac23\log2<\frac12.
 \label{p13:rig:nonaligned-bound}
\end{equation}
Compatibility of the two events is not needed.
\end{theorem}

\begin{proof}
Suppose along a sequence that
$d_{xy}\ge((2/3)\log2+\delta)B$ for some fixed
$\delta>0$.  Choose a fixed $a<1/2$, sufficiently close to
$1/2$, that
$(2/3)\log(1/a)<(2/3)\log2+\delta/2$.
Theorem~\ref{p13:rig:short-basis} gives a space $V$
of dimension at least $\delta B/2-o(B)$ with injective
physical projections and a basis whose half-cardinalities
are at most $q_0=B^{a+o(1)}$.  The quantitative bridge to
the common-curve theorem is
\begin{equation}
              q_0^2\log(B+1)=B^{2a+o(1)}=o(B)=o(\log N).
 \label{p13:rig:threshold-bridge}
\end{equation}
Thus Theorem~\ref{p13:rig:common-curve} supplies an
irreducible $C_0$ through $(x,y)$ of degree
$D\le B^{a+o(1)}$ and logarithmic height $o(B)$, with
functional image of dimension at most one.

If $D\ge2$, Proposition~\ref{p13:rig:grid-fibres} gives
at least $\delta B/2-o(B)$ integer points in the offset
box, because $D^2=o(B)$.  Corollary~\ref{p13:rig:sublinear-grid}
gives $o(B)$ points, a contradiction.

If $D=1$, write $C_0=\alpha X+\beta Y+\gamma$.
Both $\alpha,\beta$ are nonzero.  Since
$H(C_0)=\exp(o(B))=N^{o(1)}$ and $x,y\ge N$, they have
opposite signs; otherwise $\gamma$ could not cancel the
two terms.  The line therefore has the form
$X=\lambda Y+\nu$ with $\lambda>0$.
The kernel of the functional-class map has positive
dimension.  Choose a nonzero word in it.  Each selection
is nonempty and even by projection injectivity.  On the
line, the distinct simple roots from the two blocks must
pair, each pairing giving $\nu+u=\lambda v$.
There are at least two pairings, hence
\[
 \lambda=\frac{u_1-u_2}{v_1-v_2}=\frac pq,
 \qquad (p,q)=1,\qquad 1\le p,q\le2B.
\]
Evaluating one pairing at the centres gives
$qx-py=pv-qu$, of absolute value at most $(p+q)B$.
This is an alignment as in \eqref{p13:rig:alignment},
contrary to the assumption.  Both degree cases are
impossible, and $\delta$ was arbitrary.
\end{proof}

The number $(2/3)\log2$ is a sufficient upper threshold
arising from the extraction and degree bounds.  No matching
rank construction is asserted for this constant.


\subsection{Aligned pairs and the exact contact distance}

For an aligned pair the line is already known.  Substitution into it
costs a height linear in the support, so the extraction exponent may
approach one without the quadratic resultant threshold.

\begin{proposition}[Rank on a known alignment]
\label{p13:rig:aligned}
Suppose $bx-ay=h$, where $(a,b)=1$, $1\le a,b\le B^3$,
and $|h|\le(a+b)B$.  For every fixed $\tau\in(0,1)$,
\begin{equation}
 d_{xy}\le\frac B{\max(a,b)}+\frac23B\log(1/\tau)
                    +O_{\tau,C}(B/\sqrt{\log B}).
 \label{p13:rig:aligned-quantitative}
\end{equation}
For distinct compatible centres this implies, uniformly
over aligned pairs,
$d_{xy}\le B/2+o_C(B)$.  If such a pair has an alignment
with $a=b=1$, then $d_{xy}=o_C(B)$.
\end{proposition}

The proof in Section~\ref{p14:comp:rigidity:aligned} counts the matched
linear factors on that line.  At height one, compatibility imposes
$|x-y|\ge L$; distance exactly $L=B-1$ leaves only one shared true
position.  A nonzero even relation needs at least two matches, so this
contact contributes no macroscopic intersection rank.  Thus the bound
includes the minimum compatible separation exactly.

\subsection{Uniform triple nullity and its probabilistic cost}

\begin{corollary}[Triple rank and maximum joint probability]
\label{p13:rig:triple}
Every compatible triple $T$ of distinct centres in the
fixed critical band satisfies
\begin{equation}
                 \rho_{\rm rel}(T)\le\frac32L+o_C(L).
 \label{p13:rig:triple-nullity}
\end{equation}
If any of its pairs is nonaligned, the stronger bound
$\rho_{\rm rel}(T)\le(1+(2/3)\log2)L+o_C(L)$ holds.
In particular, with $q(T)=\E\prod_{x\in T}J_x$,
\begin{equation}
                  \max_{\substack{T\subset I_N\\|T|=3}}q(T)
                            \le N^{-3/2+o_C(1)}.
 \label{p13:rig:max-triple}
\end{equation}
\end{corollary}

\begin{proof}
For three spaces of dimensions at most $L$, the dimension
identity and the intersection bound give
\[
 \begin{split}
 \rho_{\rm rel}(T)
 &=\dim(S_1\cap S_2)+\dim((S_1+S_2)\cap S_3)\\
 &\le L+\dim(S_1\cap S_2).
 \end{split}
\]
Permuting the indices gives
$\rho_{\rm rel}(T)\le L+\min_{i<j}d_{ij}$.
Every pair in a compatible triple is compatible.
Theorem~\ref{p13:rig:nonaligned} and
Proposition~\ref{p13:rig:aligned}, with $B=L+1$, prove
the stated nullity bounds.

For clarity, the passage to probabilities uses exact
affine constraints.  The prime signs induce a uniform
character on each finite square-class space.  The
constraints for $J_x=1$ have classes
$\{x-1,x\}$ and $\{x,x+j\}$, $1\le j<L$, with prescribed
values $-1$ and $+1$ respectively.  These selections span
$E_x$.  Thus an admissible window has
$p_x:=\E J_x=2^{-\dim S_x}$, and a compatible family has
\begin{equation}
 q(T)=2^{-\dim\sum_{x\in T}S_x}
                  =2^{\rho_{\rm rel}(T)}\prod_{x\in T}p_x.
 \label{p13:rig:exact-probability}
\end{equation}
Every local relation is supported on defective positions,
so Lemma~\ref{p13:rig:defective} implies
$L-O_C(B/\log B)\le\dim S_x\le L$.
Consequently $p_x=2^{-L+o_C(L)}$ for admissible windows;
inadmissible windows contribute zero.  Substitution in
\eqref{p13:rig:exact-probability} gives
$q(T)\le2^{-3L/2+o_C(L)}=N^{-3/2+o_C(1)}$ uniformly,
which is \eqref{p13:rig:max-triple}.
\end{proof}

The estimate \eqref{p13:rig:max-triple} is the arithmetic
input needed by the reconstruction from three-window
activities.  Its proof uses the true windows and exact
marginals, so no removal of locally exceptional windows
is required at this stage.

\subsection{The short-basis and common-curve constructions}
\label{p14:comp:rigidity}

The notation is that of Section~\ref{p13:rig:section}.  This section
supplies the complete local estimates, basis extraction, and algebraic
construction used by the main rank argument.  The final
degree-uniform point count and the nonaligned contradiction are proved
in this section.

\subsubsection{Local square relations and defective positions}

For primes larger than the window length, a local square relation
can use only positions having no odd valuation at any such prime.
The next elementary argument bounds those positions without an
admissibility assumption.

\begin{lemma}[A local square threshold]
\label{p13:rig:local-square}
Let $B\ge2$ be an integer, let $I\subset[-B,B]\cap\Z$ have positive even
cardinality $m$, and let $x$ be a positive integer.  If
$x>(B+1)^{5m}$, then $\prod_{a\in I}(x+a)$ is not an integer
square.
\end{lemma}

\begin{proof}
The assumed threshold implies $x\ge2B$.
If the product were a square, Lemma~\ref{p1:arith:runge},
with $R=B$ and $d=m$, would give $x\le(64B)^m$.
For $B\ge2$, we have $64B\le(B+1)^5$.
Thus $x\le(B+1)^{5m}$,
contrary to the hypothesis.
\end{proof}

\begin{lemma}[Few defective positions in every critical window]
\label{p13:rig:defective}
The number of $n\in V_x$ for which
\[
             \prod_{\substack{p>10B\\v_p(n)\ {\rm odd}}}p=1
\]
is $O_C(B/\log B)$, uniformly over all critical windows.
\end{lemma}

\begin{proof}
Apply the local assertion of Lemma~\ref{p1:arith:defects}
with $H=10B$ to $V_x$.
In the critical band its diameter is $O(B)$ and its values
lie in $[N/2,3N]$ for all sufficiently large $N$.
The comparison constants are uniform in $x$.
\end{proof}

We use the elementary bound $\pi(t)\ll t/\log t$.  For
completeness, the product of the primes in $(u,2u]$ divides
$\binom{2u}{u}\le4^u$ for integer $u$; dyadic summation gives
$\vartheta(t)=\sum_{p\le t}\log p\ll t$.  Splitting at
$\sqrt t$ gives the stated bound for $\pi(t)$.

We will also need the elementary prime harmonic estimate
\begin{equation}
 \sum_{B^a<p\le10B}\frac1p
       =\log(1/a)+O_a(1/\log B),\qquad 0<a<1\text{ fixed}.
 \label{p13:rig:mertens}
\end{equation}
Here is a derivation recording its error term.  Chebyshev's bound
above implies $\psi(t)=\sum_{p^j\le t}\log p\ll t$.
The identity
$\log(n!)=\sum_{d\le n}\Lambda(d)\lfloor n/d\rfloor$
then gives $\sum_{d\le t}\Lambda(d)/d=\log t+O(1)$.
The contribution of prime powers with exponent at least two is
bounded, so $A(t)=\sum_{p\le t}(\log p)/p=\log t+O(1)$.
Partial summation yields
\[
 \sum_{p\le t}\frac1p
 =\frac{A(t)}{\log t}+\int_2^t\frac{A(u)}{u(\log u)^2}\,du
 =\log\log t+c_0+O(1/\log t).
\]
The difference at $10B$ and $B^a$ proves
\eqref{p13:rig:mertens}.

\subsubsection{Large primes turn intersections into components}

Construct a graph on the disjoint formal copies of $V_x,V_y$.
A prime $p>10B$ has at most one occurrence of odd valuation
in either block.  If it has two occurrences, join their vertices;
if it has only one, mark that vertex as forced to zero.
The corresponding binary equation says that the endpoints of
an edge have equal coordinates, or that a marked coordinate is
zero.  A component is free when it contains no marked vertex.

\begin{lemma}[Free components]
\label{p13:rig:components}
Let $c$ be the number of free components containing an edge.
Then
\[
                         c=d_{xy}+O_C(B/\log B).
\]
Each such component meets both blocks and is selected in its
entirety or omitted in any solution to the large-prime equations.
\end{lemma}

\begin{proof}
The isolated unmarked vertices are precisely the defective
positions.  Their total number $D_0$ is $O_C(B/\log B)$ by
Lemma~\ref{p13:rig:defective}.  The kernel of the large-prime
equations has dimension $c+D_0$.  Imposing the remaining
prime equations and the two block parities costs at most
$\pi(10B)+2$ dimensions.  If $\mathcal R$ is the space of
complete relations and $k_x=\dim\ker\phi_x$, then the map
taking a relation to its common class gives an exact sequence
with kernel $\ker\phi_x\oplus\ker\phi_y$.  Hence
\[
 c+D_0-\pi(10B)-2\le\dim\mathcal R=k_x+k_y+d_{xy}
                                      \le c+D_0.
\]
A local relation can use only defective positions, so
$k_x+k_y\le D_0$.  These inequalities prove the assertion.
\end{proof}


\subsubsection{Proof of Theorem~\ref{p13:rig:short-basis}}
\label{p14:comp:rigidity:short-basis}

\begin{proof}
\emph{A graph estimate.}
A multigraph with $n$ vertices and $e/n\to\infty$ has a cycle
of length $O(1+\log(n+1)/\log(e/n))$.  Loops and parallel
edges give cycles of length one and two.  Otherwise delete
vertices of degree less than $t=\lfloor e/(2n)\rfloor$.
Not all vertices can disappear, since fewer than $nt\le e/2$
edges would then be deleted.  The remaining graph has minimum
degree at least $t$.  If its girth exceeds $2r+1$, a breadth
first search contains at least
$1+t\sum_{j<r}(t-1)^j$ vertices.  This bounds $r$ and proves
the estimate.

Consequently, when $n\ll B/\log B$ and $e\le B$, the binary
cycle space contains at least $e-O(B/\sqrt{\log B})$
independent cycles of length $O(\log B/\log\log B)$.
Indeed, while at least $B/\sqrt{\log B}$ edges remain,
record a short cycle and delete just one of its edges.
The recorded cycles are independent: the first deleted edge
belongs to no later cycle, and the same argument applies
successively to all deleted edges.

\emph{Medium-prime columns.}
Start with the free nontrivial components of
Lemma~\ref{p13:rig:components}, and discard those with more
than $M=\lfloor\log B\rfloor$ vertices.  This loses at most
$2B/(M+1)$ components.  The remaining number is
$d_{xy}+O_C(B/\log B)$.  Put $z=B^a$.  To each retained
component attach its parity column on the primes $z<p\le10B$,
and let $w_j$ be that column's Hamming weight.  Each nonzero
entry consumes a divisible position, whence
\[
 \sum_j w_j
 \le\sum_{z<p\le10B}\#\{n\in V_x\sqcup V_y:p\mid n\}
 \le2B\sum_{z<p\le10B}\frac1p+2\pi(10B)
 \le2B\log(1/a)+O_a(B/\log B).
\]
Discard only columns of weight at least three.  Their number
is at most one third of this sum, so at least
\[
                 e\ge d_{xy}-\frac23B\log(1/a)
                                   -O_{a,C}(B/\log B)
\]
columns remain.  There is no factor $M$ in this count: it
bounds the number of discarded columns, not their physical size.

Use the primes $z<p\le10B$ and one extra vertex $\infty$
as vertices of a multigraph.  A column of weight two joins its
two primes; a column of weight one joins its prime to $\infty$;
a zero column is a loop at $\infty$.  Its binary incidence
matrix has the same kernel as the parity matrix, because the
row at $\infty$ is the sum of the other rows.  The graph has
$n\le\pi(10B)+1\ll B/\log B$ vertices and $e\le B$ edges.
The preceding graph argument therefore supplies a space
$\mathcal U$ of dimension at least
$e-O(B/\sqrt{\log B})$ with a basis of short cycles.
Each basis vector uses
$O(\log^2 B/\log\log B)$ physical positions.

\emph{Small primes and parity.}
All large and medium primes already cancel on $\mathcal U$.
The primes at most $z$ and the two block parities impose at
most $s=\pi(z)+2\ll_a B^a/\log B+2$ further equations on its
cycle coordinates.  Their kernel is $\mathcal K_a$; since
$s=o(B/\sqrt{\log B})$, it has the dimension in
\eqref{p13:rig:short-dimension}.  A fundamental kernel basis
for a matrix with at most $s$ rows uses at most $s+1$ columns
per vector: express each nonpivot column in terms of pivot
columns.  Thus each resulting vector uses at most $s+1$
short cycles, giving \eqref{p13:rig:short-support}.

Finally, the original free nontrivial components are disjoint
and meet both blocks.  Every nonzero selection of them contains
a component carrying a large prime which occurs once in each
block and in no other component.  Its common square class is
therefore nonzero.  This also proves injectivity of both physical
projections and of the common-class map.  The imposed parity
equations complete the proof.
\end{proof}

\subsubsection{Two relations produce a polynomial of controlled height}

The short basis will have support $B^{a+o(1)}$.  We therefore
need an auxiliary polynomial whose degree and logarithmic
height grow linearly with the support, and a simultaneous
construction whose threshold is quadratic in that support.
The following elementary factor bound will keep the height
under control when common polynomial factors are removed.

\begin{lemma}[Height of a bivariate factor]
\label{p13:rig:factor-height}
Let $P\in\Z[X,Y]\setminus\{0\}$, with partial degrees
$d_X,d_Y$.  Every primitive integer factor $Q$ of $P$, and
the corresponding integer quotient, satisfy
\[
 H(Q)\le2^{d_X+d_Y}\|P\|_2
       \le2^{d_X+d_Y}\sqrt{(d_X+1)(d_Y+1)}\,H(P).
\]
\end{lemma}

\begin{proof}
The multivariate Mahler measure $M(P)$ is multiplicative,
at most $\|P\|_2$, and at least one for a nonzero integer
polynomial.  Jensen's formula in one variable, followed by
the leading coefficient in that variable, proves the last
assertion by iteration.  The one-variable coefficient bound
also iterates: integrating first in $X$ and then in $Y$ gives
\[
 |[X^iY^j]P|\le\binom{d_X}{i}\binom{d_Y}{j}M(P).
\]
Thus $H(Q)\le2^{\deg_XQ+\deg_YQ}M(Q)$; multiplicativity
and $M(P/Q)\ge1$ prove the result.  The quotient has the same
bound, with the integer content handled separately.
\end{proof}

\begin{proposition}[Two complete witnesses]
\label{p13:rig:two-witnesses}
There is an absolute constant $c>0$ with the following property.
Let $u,v$ be independent formal words, even in each block,
with offsets in $[-B,B]$.  For $w=u,v,u+v$, write
$|I_w|=2r_w$, $|J_w|=2s_w$, and assume
$\max(r_w,s_w)\le q$, where $1\le q\le B$.
Suppose that $F_u(x,y)$ and $F_v(x,y)$ are integer squares,
with $x,y\in[N,2N)$.  If $N>(B+1)^{cq}$, there is a
nonzero $R\in\Z[X,Y]$ such that
\begin{equation}
 R(x,y)=0,\qquad
 \deg_XR,\deg_YR\le12q,\qquad \deg R\le24q,
 \qquad H(R)\le(B+1)^{cq}.
 \label{p13:rig:two-polynomial}
\end{equation}
Moreover, $R$ is the norm of a biquadratic auxiliary function
whose four polynomial coefficients have no common nonconstant
factor.
\end{proposition}

\begin{proof}
Set $F_1=F_u$, $F_2=F_v$, $G=\gcd(F_1,F_2)$, with monic
gcd, and $F_3=F_{u+v}=F_1F_2/G^2$.
Formal independence makes
\[
 K=\mathbb Q(X,Y)(W_1,W_2),\qquad W_i^2=F_i,
\]
a degree-four extension.  With $W_0=1$ and
$W_3=W_1W_2/G$, the four $W_\chi$ form a basis.  At the
given integer point use positive square roots.  Then all
$W_\chi(x,y)$ are integers: $F_3(x,y)$ is an integer which
is a rational square, hence an integer square, and
$W_3(x,y)=\sqrt{F_1(x,y)}\sqrt{F_2(x,y)}/G(x,y)$.

\emph{The linear system at infinity.}
Put $r_0=s_0=0$, $k=3q$, and consider
\[
 \mathcal A=\sum_{\chi=0}^3A_\chi(X,Y)W_\chi,
 \qquad \deg_XA_\chi\le k-r_\chi,
 \quad \deg_YA_\chi\le k-s_\chi.
\]
There are
\[
 n=\sum_{\chi=0}^3(k-r_\chi+1)(k-s_\chi+1)
   \ge21q^2+18q+4
\]
unknown coefficients, and $n\le4(3q+1)^2$.
Writing $\xi=1/X$, $\eta=1/Y$, the positive formal branches
are
\[
 W_\chi=X^{r_\chi}Y^{s_\chi}
 \prod_{u\in I_\chi}(1+u\xi)^{1/2}
 \prod_{v\in J_\chi}(1+v\eta)^{1/2}.
\]
Thus $f(\xi,\eta)=\xi^k\eta^k\mathcal A$ is a formal
power series.  Require its coefficients of total degree at
most $2k$ to vanish.  This gives at most
\[
 t=\frac{(2k+1)(2k+2)}2=18q^2+9q+1,
 \qquad n-t\ge3q^2+9q+3,\qquad \frac t{n-t}<6
\]
homogeneous linear equations.

All coefficients required from the binomial series have
denominator dividing $4^{2k}$.  Cauchy's estimate on circles
of radius $(2B)^{-1}$ bounds their numerators, after this
common denominator is cleared, by $(B+1)^{c_0q}$ for an
absolute $c_0$.  An elementary pigeonhole form of Siegel's
lemma yields a nonzero integer solution of height
$(B+1)^{c_1q}$.  Indeed, if the integer matrix has height $A$,
the images of $\{0,\ldots,M\}^n$ take at most
$(2nAM+1)^t$ values.  Choosing
$M=\lceil(3nA)^7\rceil$ makes two images equal, because
$t/(n-t)<6$.  Their difference is a nonzero kernel vector;
$\log n=O(\log B)$ is absorbed in the stated height.

\emph{Primitive coefficients.}
Let $D$ be the primitive polynomial gcd of the four
$A_\chi$, omitting zero coefficients, and write
$A_\chi=DA'_\chi$.  The factor bound
of Lemma~\ref{p13:rig:factor-height} keeps
$H(A'_\chi)\le(B+1)^{O(q)}$.  Write
$a=\deg_XD$, $b=\deg_YD$, $e=\deg D$.  Then
\[
 f=\widetilde D f',\qquad
 \widetilde D=\xi^a\eta^bD(\xi^{-1},\eta^{-1}),\qquad
 f'=\xi^{k-a}\eta^{k-b}\sum_\chi A'_\chi W_\chi.
\]
These are power series, and
$\operatorname{ord}\widetilde D=a+b-e$ for total order.
Since order is additive in the power-series domain,
\begin{equation}
 \operatorname{ord}f'\ge2k+1-a-b+e
                          =(k-a)+(k-b)+e+1.
 \label{p13:rig:primitive-order}
\end{equation}
Thus removing $D$ preserves vanishing beyond the remaining
pole weight.  This is a formal division, never a division
by the possibly zero value $D(x,y)$.  Replace the auxiliary
by its primitive version from now on.

\emph{Integer separation and the norm.}
On the bidisc $|\xi|,|\eta|\le(4B)^{-1}$, the analytic
numerator is bounded by $(B+1)^{O(q)}$.  Cauchy's estimate
and the order of vanishing give, in the comparable range
$x,y\in[N,2N)$,
\begin{equation}
                  |\mathcal A(x,y)|\le
                                   \frac{(B+1)^{c_2q}}N.
 \label{p13:rig:integer-gap}
\end{equation}
To see the power of $N$ explicitly before removing $D$, the
tail of the numerator is bounded by its coefficient bound
times
\[
 \sum_{\ell\ge2k+1}(\ell+1)(4B/N)^\ell
       \ll(k+1)(4B/N)^{2k+1}.
\]
Multiplication by $x^ky^k\le(2N)^{2k}$ gives
\eqref{p13:rig:integer-gap}.  After removing $D$, both pole
weights decrease, and \eqref{p13:rig:primitive-order} gives
the same estimate with the new height bound.

The specialization $\mathcal A(x,y)$ is an integer, so it
vanishes for the stated threshold.  The auxiliary is nonzero
in $K$ by independence of the four characters.  Its norm is
therefore nonzero.  An explicit polynomial expression is
\[
 \begin{split}
 U&=A_0^2+A_1^2F_1-A_2^2F_2-A_3^2F_3,\\
 V&=A_0A_1-A_2A_3(F_2/G),\\
 R&=U^2-4F_1V^2.
 \end{split}
\]
Here $F_2/G$ is a polynomial.  The partial degrees of $U$
are at most $2k$, and those of $V$ are at most
$2k-r_1$, $2k-s_1$.  For the second term in $V$, use
$r_3=r_1+r_2-\deg_XG$ and its $Y$ analogue.
It follows that the partial degrees of $R$ are at most
$4k=12q$.  The coefficient $\ell^1$ norms of all factors
in this formula are $(B+1)^{O(q)}$, which proves the height
bound and completes \eqref{p13:rig:two-polynomial}.
\end{proof}


\subsubsection{Proof of Theorem~\ref{p13:rig:common-curve}}
\label{p14:comp:rigidity:common-curve}

\begin{proof}
For $i<j$, the half-cardinalities of $w_i,w_j,w_i+w_j$
are at most $q=\min(2q_0,B)$, since an even subset of
$[-B,B]\cap\Z$ has size at most $2B$.
Proposition~\ref{p13:rig:two-witnesses} gives a nonzero
polynomial $R_{ij}$ of degree at most $48q_0$ and height
$(B+1)^{O(q_0)}$, vanishing at $(x,y)$ and arising from
an auxiliary with primitive coefficients.
Choose an irreducible primitive factor $C_0$ of $R_{12}$
which vanishes at this point.  The factor-height lemma
preserves its height bound.  If $C_0$ depended on only one
variable, its integer root of size at least $N$ would
contradict Cauchy's root bound.  Thus it depends on both.

If $C_0$ failed to divide some $R_{ij}$, the resultant
$\operatorname{Res}_Y(C_0,R_{ij})$ would be a nonzero integer
polynomial vanishing at $x$.  Its height is
$(B+1)^{O(q_0^2)}$: its Sylvester matrix has size $O(q_0)$,
entries of coefficient norm $(B+1)^{O(q_0)}$, and
$O(q_0)!$ determinant terms.  This again contradicts the
root bound under the stated threshold.  Hence $C_0$ divides
every $R_{ij}$.

No factor $X+u$ or $Y+v$ vanishes identically on this curve,
since $C_0(x,y)=0$ with both coordinates larger than $B$.
The displayed square classes are therefore defined in
$K_0$.  If the classes of $F_{w_i},F_{w_j}$ were independent
there, they would still generate a degree-four field
extension.  Specializing the polynomial norm formula gives
zero, since $C_0\mid R_{ij}$; in a field this forces the
auxiliary to vanish.  Independence of its four characters
would force all four coefficient polynomials to be divisible
by $C_0$, contrary to their primitivity.  Thus the images
of every pair of basis vectors are dependent.  Over $\F_2$
this means that all nonzero images are equal, so the image
has dimension at most one.
\end{proof}

\subsubsection{Proof of Proposition~\ref{p13:rig:aligned}}
\label{p14:comp:rigidity:aligned}

\begin{proof}
Fix $\tau<1$.  Theorem~\ref{p13:rig:short-basis} gives a
space $V=\mathcal K_\tau$ of dimension
\[
 t\ge d_{xy}-\frac23B\log(1/\tau)
                     -O_{\tau,C}(B/\sqrt{\log B}),
\]
with a basis of half-cardinalities at most
$q_0=B^{\tau+o(1)}$.  The case $t<2$ is covered by
the displayed error term.  For $t\ge2$, apply
Proposition~\ref{p13:rig:two-witnesses} to each pair of
basis vectors.  Its threshold holds because
$q_0\log(B+1)=o(B)=o(\log N)$.

Put $\ell(X,Y)=bX-aY-h$.  Every resulting norm polynomial
$R_{ij}$ is divisible by $\ell$.  Otherwise, with
$d_Y=\deg_YR_{ij}$, the nonzero integer polynomial
\[
           a^{d_Y}R_{ij}\left(X,\frac{bX-h}{a}\right)
\]
would vanish at $x$.  Its height is $(B+1)^{O(q_0)}$
since $\max(a,b,|h|)\le2B^4$, contradicting the
integer root bound for $x\ge N$.
The primitive-norm argument in
Theorem~\ref{p13:rig:common-curve} now applies to this
line and shows that the image of $V$ in
$\mathbb Q(\ell=0)^\times/\mathbb Q(\ell=0)^{\times2}$
has dimension at most one.

Let $V_0$ be its kernel, of dimension at least $t-1$.
On the line, $X+u$ and $Y+v$ have the same root precisely when
\begin{equation}
                            h+bu=av.
 \label{p13:rig:line-matching}
\end{equation}
Every word of $V_0$ must pair all its first-block roots
with second-block roots.  The integer solutions of
\eqref{p13:rig:line-matching} advance by $(a,b)$ in two
intervals of $B$ consecutive offsets, so their number $r$
satisfies
\[
                        r\le1+\frac{B-1}{\max(a,b)}.
\]
Injectivity of the physical projection gives
$t-1\le\dim V_0\le r$, proving
\eqref{p13:rig:aligned-quantitative} after absorbing the
bounded terms.

For compatible distinct centres with $\max(a,b)\ge2$,
this gives $d_{xy}\le B/2+(2/3)B\log(1/\tau)+o_{\tau,C}(B)$.
Choose $\tau$ sufficiently close to one after fixing a
desired error margin, then let $N$ tend to infinity.
This order of choices proves the uniform $B/2+o_C(B)$ bound.

For $a=b=1$, compatibility implies $|x-y|\ge L$:
if $x<y<x+L$, the first run makes $f(y-1)=f(y)$,
whereas the second requires their signs to be opposite.
At distance exactly $L=B-1$, the true windows share one
position; at larger distance they are disjoint.  Thus
\eqref{p13:rig:line-matching} has at most one solution.
A nonzero word in $V_0$ would need at least two pairings,
because its selections are nonempty and even.  Consequently
$V_0=0$, $t\le1$, and the dimension estimate followed by
$\tau\uparrow1$ gives $d_{xy}=o_C(B)$.
\end{proof}
